% capitulos/anexo.tex

\section{Critérios Planejados para a Coleta do Dataset}

A coleta de dados a ser realizada no TCC2 seguirá os critérios definidos abaixo, estabelecidos para garantir a qualidade, representatividade e rastreabilidade do \textit{dataset}:

\begin{itemize}
    \item \textbf{Volume total previsto:} 150 vídeos, distribuídos em 3 classes (50 por classe: ``biblioteca'', ``prova'' e ``secretaria'')
    \item \textbf{Participantes:} 4 voluntários surdos fluentes em Libras, em ambiente controlado
    \item \textbf{Duração por sinal:} aproximadamente 1 segundo por sequência (30 quadros a 30 FPS)
    \item \textbf{Resolução de captura:} 640$\times$480 pixels via webcam convencional
    \item \textbf{Condições de iluminação:} natural e artificial, com pelo menos uma sessão em cada condição por voluntário
    \item \textbf{Posicionamento da câmera:} frontal, enquadrando o corpo do voluntário da cintura para cima
    \item \textbf{Variações planejadas:} diferentes velocidades de execução, pequenas variações de ângulo e distância da câmera
    \item \textbf{Anonimização:} apenas os vetores de \textit{landmarks} extraídos serão armazenados; os vídeos originais não serão retidos no repositório do projeto
\end{itemize}

\section{Protocolo de Consentimento dos Participantes}

Em conformidade com os princípios éticos para pesquisa envolvendo seres humanos estabelecidos pela Resolução CNS 510/2016, todos os voluntários participantes da etapa de coleta de dados deverão assinar o Termo de Consentimento Livre e Esclarecido (TCLE) antes do início das gravações.

O TCLE informará aos participantes:

\begin{itemize}
    \item Os objetivos da pesquisa e o uso previsto dos dados coletados
    \item Que a participação é voluntária e pode ser interrompida a qualquer momento sem ônus
    \item Que os vídeos não serão armazenados após a extração dos \textit{landmarks}
    \item Que os dados serão utilizados exclusivamente para fins acadêmicos, com garantia de anonimato
    \item Os contatos dos pesquisadores responsáveis para esclarecimentos
\end{itemize}

O modelo do TCLE será elaborado em conformidade com as diretrizes do Comitê de Ética em Pesquisa (CEP) da instituição e submetido à aprovação antes do início da coleta.

\section{Requisitos Mínimos Previstos para Execução do Sistema}

Os requisitos abaixo foram definidos com base nas especificações de hardware utilizadas pelos sistemas de referência identificados na revisão bibliográfica, especialmente \citeonline{attili:24} e \citeonline{kamble:25}, que demonstraram viabilidade de execução em hardware convencional:

\begin{table}[htb]
    \centering
    \caption{Requisitos mínimos previstos para execução do LibrasVision}
    \label{tab:requisitos_sistema}
    \begin{tabular}{|l|l|}
        \hline
        \textbf{Componente} & \textbf{Especificação Mínima} \\
        \hline
        Processador & Intel Core i5 (6ª geração ou superior) ou equivalente \\
        Memória RAM & 8 GB \\
        Câmera & Webcam HD 720p (USB ou integrada) \\
        Armazenamento & 2 GB livres para modelo e dependências \\
        Python & Versão 3.8 ou superior \\
        Sistema Operacional & Windows 10+, Ubuntu 20.04+ ou macOS 11+ \\
        Conexão à internet & Necessária apenas para instalação inicial das dependências \\
        \hline
    \end{tabular}
\end{table}

O sistema não requer GPU dedicada para inferência em tempo real, o que é consistente com o objetivo de viabilizar sua execução em laboratórios e salas de aula com infraestrutura padrão.
